% Chapitre 5 : La Normalisation

\section{Chapitre 5 : La Normalisation}



Les normes sont des apports documentés contenant des spécifications techniques ou autres critères précis destinés à être utilisés systématiquement en tant que règles, lignes directrices ou définitions des caractéristiques pour assurer que des matériaux, produits, processus et services sont aptes à leurs emplois.



\subsection{La Norme ISO et la Certification ISO 9000}



L'International Organization for Standardization (ISO), créée en 1947, est une fédération mondiale d'organismes nationaux de normalisation. La norme ISO 9000 traduit la démarche qualité au niveau de l'entreprise et apporte : en interne, l'amélioration de la compétitivité et la diminution des coûts ; à l'extérieur, la réponse aux exigences des donneurs d'ordre.



Pour mettre en Å“uvre une norme ISO 9000, il faut réaliser un diagnostic de l'organisation en auditant les fonctions principales. Les points nécessaires incluent : la gestion documentaire, la revue du contrat, la gestion des produits non-conformes, et la traçabilité. Dans le contexte de SaaS Directory, la gestion documentaire est assurée par les fichiers README.md, AGENTS.md, et les conventions de commit Git. La traçabilité est garantie par l'historique Git et les migrations Drizzle numérotées (0000 à 0003).



\subsection{Les Normes IEEE}



L'Institute of Electrical and Electronics Engineers (IEEE) a défini un ensemble de 37 normes s'intéressant aux processus, produits et techniques de base du génie logiciel. Trois normes sont particulièrement pertinentes pour notre projet.



\begin{definitionbox}[title={IEEE Std 12207.0}]

Fixe la signification des termes et choisit les processus prioritaires du cycle de vie du logiciel. Peut être considérée comme un dictionnaire et une liste de vérification pour la gestion des projets.

\end{definitionbox}



\begin{definitionbox}[title={IEEE Std 830}]

Recommandée pour les spécifications des exigences du logiciel. Pourrait être intégrée avec les approches de validation par prototype.

\end{definitionbox}



\begin{definitionbox}[title={IEEE Std 1028}]

Pour appliquer la validation et la vérification. Norme facilement adaptable à tous les problèmes.

\end{definitionbox}



L'application de ces normes à SaaS Directory se concrétise par : la définition claire des exigences fonctionnelles (authentification, marketplace, messagerie, monétisation) conformément à l'IEEE 830 ; la revue systématique du code via les pull requests Git, alignée avec l'IEEE 1028 ; et la traçabilité des exigences à travers les tests end-to-end Playwright, en accord avec l'IEEE 12207.0.



\subsection{Normalisation Appliquée au Projet}



\subsubsection{Normes de Codage}



Le projet respecte des conventions strictes de normalisation du code source. Les normes de nommage sont systématiquement appliquées : PascalCase pour les composants React (\texttt{UserCard.tsx}, \texttt{ProductForm.tsx}), camelCase pour les fonctions et variables (\texttt{sendMessage()}, \texttt{upvoteCount}), et kebab-case pour les routes API (\texttt{/api/messages/send}). Les imports sont organisés par groupe : React/Next.js, bibliothèques tierces, composants internes, utilitaires. Cette normalisation facilite la revue de code et réduit les conflits de fusion.



\subsubsection{Architecture Normalisée}



L'architecture suit le pattern MVC (Model-View-Controller) adapté à Next.js 16.2.4. Le \textbf{Modèle} est constitué par le schéma Drizzle ORM (\texttt{src/lib/db/schema.ts}) et les fonctions d'accès aux données. La \textbf{Vue} est implémentée par les composants React et les Server Components de Next.js. Le \textbf{Contrôleur} est assuré par les Server Actions (\texttt{src/lib/actions/}) qui encapsulent la logique métier. Cette séparation claire respecte les principes SOLID et facilite les tests unitaires.



\subsubsection{Gestion Documentaire et Traçabilité}



La gestion documentaire est assurée par plusieurs artefacts. Le fichier README.md présente le projet, son installation et son utilisation. Le fichier AGENTS.md documente les conventions spécifiques pour les agents de développement. Le fichier package.json liste l'ensemble des dépendances avec leurs versions exactes (Next.js 16.2.4, React 19.2.4, Drizzle ORM 0.45.2, etc.). Les commits Git suivent une convention descriptive (\texttt{feat:}, \texttt{fix:}, \texttt{chore:}). Les migrations Drizzle sont numérotées séquentiellement (0000 à 0003) et documentées, garantissant la traçabilité complète de l'évolution du schéma de base de données.

